%
% sddspcas documentation.
%
\begin{sddsprog}{sddspcas}
\item \textbf{description:}
\verb+sddspcas+ is a Portable Channel Access Server (PCAS) that creates and serves EPICS process variables (PVs)
from one or more SDDS input files. Each input file must contain a \verb+ControlName+ column with PV base names.
Optional columns (e.g., \verb+Hopr+, \verb+Lopr+, \verb+ReadbackUnits+, \verb+ElementCount+, \verb+Type+) control
limits, engineering units, waveform length, and data type.

With EPICS Base 7, PCAS is normally a separate EPICS module. The SDDS-EPICS
build enables \verb+sddspcas+ only when it finds \verb+casdef.h+ and the PCAS
libraries in EPICS Base or in a sibling \verb+pcas+ module. All other toolkit
programs can be built without PCAS.

If an \verb+Equation+ column is present, PV values can be automatically updated based on expressions involving
other PVs.

SDDS6 files may contain additional enum metadata. The PV configuration columns
retain their documented primitive types: SDDS enum values in these columns,
including limits and element counts, are rejected. Use \verb+sddscast+ for
explicit conversion. EPICS enum PVs still use string \verb+Type+ and
\verb+EnumStrings+ columns.

On Unix-like systems, if a second \verb+sddspcas+ instance is launched under the same user account while a first instance
is already running on the same host, the second instance will forward its PV definitions to the running instance and
exit immediately. The running instance will create any PVs it is not already hosting and (unless started with
\verb+-standalone+) will append the new PVs to the configured \verb+-pcasPVFile+.

\item \textbf{examples:}
\begin{verbatim}
# Serve PVs described by an SDDS file
sddspcas pvlist.sdds

# Add a prefix to all served PVs and run for 10 minutes
sddspcas pvlist.sdds -pvPrefix=TEST: -executionTime=600

# Enable periodic noise updates
sddspcas pvlist.sdds -pvPrefix=TEST: -noise=1

# Run without checking master PV lists
sddspcas pvlist.sdds -standalone
\end{verbatim}

\item \textbf{synopsis:}
\begin{verbatim}
usage: sddspcas <inputfiles>
      [-masterPVFile=<filename>]
      [-pcasPVFile=<filename>]
      [-debugLevel=<debug level>]
      [-executionTime=<seconds>] (default=12 hours, 0=forever)
      [-pvPrefix=<PV name prefix>]
      [-noise=<rateSeconds>]
      [-arraySize=<bytes>]
      [-runControlPV=string=<string>,pingTimeout=<value>]
      [-runControlDescription=string=<string>]
      [-standalone]
\end{verbatim}

\item \textbf{switches:}
\begin{itemize}
  \item {\tt -masterPVFile} --- SDDS file containing IOC PV names. Used to ensure no duplicate PVs are created.
  \item {\tt -pcasPVFile} --- SDDS file tracking PVs created by \verb+sddspcas+. Used to prevent multiple servers
        from advertising the same PVs and to record the serving host.
  \item {\tt -standalone} --- Do not use the master PV files unless explicitly provided.
  \item {\tt -executionTime} --- How long to run in seconds. Values less than 0.5 are treated as ``forever''.
  \item {\tt -pvPrefix} --- Prefix prepended to all PV names.
  \item {\tt -noise} --- Periodically adds small random noise to numeric PV values.
  \item {\tt -arraySize} --- Sets \verb+EPICS_CA_MAX_ARRAY_BYTES+ for waveform transfers.
  \item {\tt -debugLevel} --- Increase diagnostic output.
  \item {\tt -runControlPV} / {\tt -runControlDescription} --- Optional run-control integration (if built with runcontrol).
  \item {\tt -aliasCount}, {\tt -syncScan} --- legacy option names that are
        recognized but not supported by this implementation; either option
        terminates the program with an error.
\end{itemize}

\item \textbf{see also:}
\begin{itemize}
  \item \progref{sddsSoftIOC}
  \item \progref{cavget}
  \item \progref{cavput}
\end{itemize}

\item \textbf{author:} Robert Soliday, ANL/APS.
\end{sddsprog}
